% !TeX root = ../../main.tex
% courb/errata.tex -- one table. Extraction provenance and verification: tables/README.md
% \small is set INSIDE \begin{table} here, and OUTSIDE the environment in the two longtable
% errata files (tables/frame/bib_errata.tex, tables/cbic/errata.tex, each wrapped in an
% explicit {\small ... } brace pair). Both forms are correct and they are not interchangeable:
% \begin{table} is a float and opens a group, so a font change inside it ends with the table,
% whereas longtable opens no group and needs the braces to stop \small leaking into the rest of
% the document. Do not harmonize the two families. The full reasoning and the six commits the
% missing brace cost are in the header of tables/frame/bib_errata.tex.
\begin{table}[htb]
\caption{Content errata of the CoUrb 2026 article and the corrections applied in
Chapter~\ref{ch:courb} (English translation of the published text).}
\label{tab:apx:courb-errata}
\centering
\small
\begin{tabular}{@{}p{0.42\textwidth}p{0.52\textwidth}@{}}
\toprule
\textbf{Defect in the published article} & \textbf{Correction in the chapter} \\
\midrule
Results and conclusion state wins over the baseline in ``16 of the 21
evaluated combinations''. &
Audited count: 15 of the 21 combinations, with one additional technical tie
(\emph{Outdoors} in Florida, where the baseline mean exceeds the best variant by
0.02 percentage points, within one standard deviation). \\
\addlinespace
Introduction, results, and conclusion state ``average gains of 20 to 24
percentage points'' on the category task. &
Audited per-state means: gains of 20.2 to 22.0 percentage points, considering the
better of the two spatial encoders in each combination (the disclosure the audit
requires). \\
\addlinespace
Introduction states ``wins in most scenarios''. &
``outperforms the baseline in most scenarios'' (verdict-verb rule; claim
strength unchanged). \\
\addlinespace
The methodology states that maximizing the product of the task utilities ``ensures that the
update is beneficial for all tasks simultaneously''. The guarantee is unconditional as
written. &
Narrowed to the guarantee the method's own paper derives: at points that are not Pareto
stationary, and under its assumption that the task gradients there are linearly independent,
the update direction is a descent direction for every task. The paper states the property in
those terms, and it approximates the bargaining solution by an iterative procedure rather than
solving it exactly. The claim is weaker than the published one, and no result depends on
it. \\
\addlinespace
The methodology attributes to a next-POI recommendation method the claim that cyclical
mobility regularities carry discriminative information about the functional nature of the
visited places. That method models long- and short-term user preference and states nothing
about temporal signal revealing what kind of place a visit is. &
Split into the part the literature supports and the part it does not. The regularity itself is
attributed to the paper that published the check-in collection this chapter uses, already
cited elsewhere in the chapter, and the recommendation method is kept for the property it does
establish, that models of the next visit read temporal structure explicitly. The claim about
place function is withdrawn rather than re-attributed, since no work in the bibliography
supports it. The paragraph's argument, that the original architecture lacks an explicit
temporal representation, is unchanged. \\
\addlinespace
The category-encoder loss attributes its hierarchical regularization term, an L2 penalty that
pulls a subcategory embedding toward its parent over a known label tree, to a paper on
nonlinear dimensionality reduction. The two share a graph Laplacian and nothing else. &
Re-attributed to a multitask embedding method for place-category annotation, which
regularizes relatedness among categories using a predefined category hierarchy: the same
construction the loss implements. The reference was already in the bibliography and already
cited in this chapter. No term, coefficient, or reported value changes. \\
\addlinespace
Two sentences of the methodology name a data column by its identifier in the source files,
\texttt{fclass}, which is not a term a reader of the chapter can resolve. &
Replaced by the registered reader-facing name, fine class, in both sentences, set in roman
because it is ordinary English rather than a code token. The parenthetical gloss that
introduced the column is kept, so the sentence still tells the reader what a secondary
category is. \\
\addlinespace
% [2026-09-02, emphasis] Same ruling as the CBIC emphasis row. Here the published tables already
% bolded the row maximum; what was missing was the second-best mark and any statement of the rule.
% Applied by the same script and the same digit-preservation proof.
Both results tables mark the largest value in each row in bold and mark nothing else, and
neither caption states what the mark means. &
The second largest value in each row is now underlined and both captions state the rule, so
the two tables read under the same convention as the rest of the document. No value
changes; only which cells carry emphasis, and the captions. \\
\addlinespace
% [2026-09-02, B-20] The conclusion said "consistently outperforms"; the very next paragraph of the
% same section reports 15 of 21 combinations plus one technical tie, and the results section says the
% baseline retains six. Those counts are themselves an errata-round correction of the published
% "16/21 (76%)" (NORTH_STAR §4 Ch.4), so the frame had already lowered the number and left the
% adverb. Author ruling 2026-09-02: "concordo podemos trocar".
The conclusion states that the proposed model \emph{consistently} outperforms the
baseline. The paragraph that follows reports fifteen of twenty-one evaluated combinations
with one further technical tie, and the results section states that the baseline retains
six of them. &
The adverb is replaced by the scope the counts support: the model outperforms the baseline
in most of the evaluated combinations. The counts in the following paragraph are unchanged
and carry the exact figures. \\
\addlinespace
% [2026-09-02, B-22] The reader asked "o que e dependente entre modelo e dataset?" -- the sentence
% told him to read the results with caution without saying WHAT ages with the dataset. The author
% ruled to change the published chapter text and register the errata, without pointing at it from the
% text: "podemos mudar o texto do capitulo para ficar melhor e criar uma errata, mas nao precisa
% referenciar a errata". That is the deposited document's own policy -- it cites only itself and the
% repository. ⚠ NO FORWARD POINTER to Chapter 5 was written inside the reproduced conclusion: the
% empirical answer exists there (Istanbul, a non-U.S. dataset) but a cross-chapter reference inside an
% article body is dissertation structure in reproduced text, which this project does not do. The
% chapter preface carries the bridging role.
The conclusion names the exclusive use of a single dataset as a limitation and asks the
reader to interpret the generalization with caution, without saying which properties of
the results depend on that dataset. The single national context the corpus covers is not
named at all. &
The sentence now names both dependencies: the places and visiting habits the corpus
records, which belong to its period, and the single country it covers. The verdict is
scoped accordingly, as evidence about that period and that country whose generalization to
current patterns, or to cities elsewhere, remains to be tested. No result changes. \\
\bottomrule
\end{tabular}
\end{table}
